\section{Uncertainty Estimation for SAR Level-1 Products} \label{sec:level1}

This section applies the principles of the \ac{GUM} to \ac{SAR} Level-1
products by taking the \emph{product}---rather than an individual calibration
activity---as the measurand-oriented starting point. This is consistent with
metrological practice: uncertainty analysis should begin from the reported
quantity, define its measurement model, identify all relevant input quantities,
and then propagate the corresponding uncertainties to the final product
statement \cite{JCGMGUM,QA4EO3}. For Level-1 \ac{SAR}, this means that the
uncertainty analysis should be organized around the calibrated single-look complex (SLC), the
polarimetric single-look complex, the interferometric product, geocoded Analysis Ready Data
such as Normalised Radar Backscatter, and the geophysical
Doppler product.

Radiometric, geometric, and polarimetric calibration are therefore treated here
as \emph{contributors} to the uncertainty budget of the corresponding Level-1
product, not as separate top-level measurands. For example, a calibrated
single-look complex (SLC) product inherits uncertainty contributions from radiometric gain and
noise corrections, geometric timing and orbit knowledge, processor
approximations, and, where relevant, polarimetric channel corrections. In the
same way, an interferometric product inherits uncertainties from the two input
single-look complex (SLC) products, their co-registration, the baseline knowledge, and the
statistics of the interferometric estimator. An NRB/ARD product inherits the uncertainty of the
source SAR acquisition and adds uncertainty from radiometric terrain correction, DEM and ancillary
data, geocoding, resampling, masks, and map-grid definition.

The following product-based structure is adopted:
\begin{itemize}
  \item a general GUM-oriented framework for Level-1 products;
  \item a calibrated single-channel single-look complex (SLC) product as the main Level-1
  example;
  \item a polarimetric single-look complex (SLC) product, for which the scattering matrix is the
  principal measurand;
  \item an interferometric product, whose principal observables are coherence,
  phase, and related quality indicators;
  \item a geocoded ARD/NRB product, for which the measurand is terrain-normalised
  backscatter on a defined map grid;
  \item the geophysical Doppler product, treated as a detailed worked example.
\end{itemize}

% \section{Applicability of Metrology Standards to SAR Cal/Val procedures}

\subsection{General GUM approach for Level-1 SAR products}

For a generic Level-1 product measurand $y$, the measurement model can be
written in compact form as
\begin{equation}
  y = f\!\left(\mathbf{x}_{\rm scene}, \mathbf{x}_{\rm inst},
  \mathbf{x}_{\rm proc}, \mathbf{x}_{\rm aux}\right) + \Delta y,
  \label{eq:l1_product_generic_model}
\end{equation}
where $\mathbf{x}_{\rm scene}$ contains the scene-dependent quantities,
$\mathbf{x}_{\rm inst}$ the instrument and calibration quantities,
$\mathbf{x}_{\rm proc}$ the processor-dependent quantities, and
$\mathbf{x}_{\rm aux}$ the auxiliary data, such as orbit, attitude,
propagation, and reference target information. The residual term $\Delta y$
collects unmodelled effects.

Following the \ac{GUM}, the combined standard uncertainty is obtained either
through the law of propagation of uncertainty or, when the model is strongly
nonlinear, through Monte Carlo propagation as discussed in
Section~\ref{sec:uncertainty_propagation}. In covariance form,
\begin{equation}
  u^{2}(y) \approx \mathbf{J}_{f}\, \mathbf{U}_{\mathbf{x}}\,
  \mathbf{J}_{f}^{\mathsf{T}},
  \label{eq:l1_product_generic_cov}
\end{equation}
where $\mathbf{J}_{f}$ is the Jacobian of the measurement model and
$\mathbf{U}_{\mathbf{x}}$ is the covariance matrix of the input quantities.
For Level-1 \ac{SAR} products, the practical difficulty is not the propagation
itself, but the correct identification of structured errors and spatial or temporal correlations.

Across all Level-1 products, the dominant uncertainty sources generally belong
to the following categories:
\begin{itemize}
  \item instrument and internal calibration terms, such as gain drifts,
  receiver noise, antenna pattern uncertainty, timing offsets, and channel
  imbalance;
  \item focusing and processor terms, such as resampling, interpolation,
  co-registration, annotation approximations, and finite-resolution effects;
  \item scene-dependent terms, such as speckle, decorrelation,
  and target-background contrast;
  \item auxiliary-data terms, such as orbit, attitude, digital elevation model (DEM), and surveyed reference-target uncertainties.
\end{itemize}


The uncertainty budget shall also include the Cal/Val reference and the
measurement tool used to estimate the product quality indicator. For point
target IRF, geolocation, and radiometric checks, this includes target survey
uncertainty, target calibration uncertainty, signal-to-clutter ratio,
background estimation, atmospheric and ionospheric propagation at the time of
acquisition, and the implementation uncertainty of the analysis software. For
distributed-targets, it includes the spatial and
temporal collocation thresholds, representativeness error, model assumptions,
and the covariance between the satellite product, reference data, and any
background or model field.

Auxiliary data files are part of the measurement model. Orbit and attitude products, DEMs, meteorological fields, ionospheric maps,
antenna patterns, calibration look-up tables, masks, and processor annotation
files shall therefore be assigned uncertainty terms, validity domains, and
correlation scales. Atmospheric propagation terms include atmospheric
path delay and attenuation; ionospheric terms include dispersive group delay, and Faraday rotation. These
terms can affect SLC, and interferometric products and are
not limited to geocoded products.

Per-pixel uncertainty reporting is required where the product provides a
spatially varying measurand. A per-pixel uncertainty layer should contain the
local random component of the standard uncertainty, while product-level
metadata should document common-mode and spatially correlated terms that cannot
be interpreted as independent pixel noise. The recommended reporting structure
is therefore: per-pixel random uncertainty, local-window uncertainty for
estimator or multilooking effects, swath-scale or burst-scale uncertainty for
calibration and processing terms, and product-wide uncertainty for reference,
orbit, processor-version.

For consistency in the traceability diagrams used below, $u(x)$ denotes the
standard uncertainty of input quantity $x$, $\Delta x$ denotes a residual or
model error term, $\partial y/\partial x$ denotes a sensitivity coefficient,
$\mathbf{U}$ denotes a covariance matrix, $\mathbf{J}$ denotes a Jacobian, and
arrows indicate propagation from an input quantity or uncertainty contribution
to the measurand.


This product-based organization places the role of radiometric,
geometric, or polarimetric calibration inside the measurement model of the final product.

\subsection{Calibrated Level-1 SLC product}

A calibrated Level-1 single-look complex (SLC) is the natural reference example because it is the product from which most other Level-1 quality indicators are derived. Its measurand is the calibrated complex pixel value together with the associated radiometric and geometric interpretation.

\subsubsection{Step 1: Definition of measurand and measurement model}

Let the calibrated complex sample be denoted by $s_{\rm c}$. A generic model is
\begin{equation}
  s_{\rm c} = C_{\rm rad}\, C_{\rm geo}\,
  \left(s_{\rm raw} - n\right) e^{-j\phi_{\rm corr}} + \Delta s,
  \label{eq:slc_model}
\end{equation}
where $s_{\rm raw}$ is the focused but not yet fully corrected complex sample,
$n$ is the thermal-noise contribution, $C_{\rm rad}$ summarizes radiometric
calibration factors, $C_{\rm geo}$ summarizes geometric and timing corrections,
$\phi_{\rm corr}$ collects phase corrections applied during processing, and
$\Delta s$ represents residual processor and modelling errors.

Equation~\ref{eq:slc_model} makes clear that radiometric and geometric
calibration are not separate measurands here; they are input quantities in the
measurement model of the Level-1 single-look complex (SLC) product.



\subsubsection{Step 2: Traceability diagram for the single-look complex product}

Figure~\ref{fig:slc_traceability} summarizes the uncertainty propagation from
the calibrated single-look complex sample to three branches of Level-1
observables. For a calibrated single-look complex (SLC), the uncertainty budget is conveniently organized into three observable-oriented branches, each derived from the same calibrated complex sample but leading to a different product-level uncertainty statement.

The radiometric branch includes uncertainty of the absolute calibration factor,
receiver gain, antenna pattern normalization, thermal-noise subtraction,
scalloping corrections, and finite-look image statistics. When image intensity
or $\sigma^{0}$ is estimated from the single-look complex (SLC), these terms dominate the radiometric quality statement.

The geometric branch includes range timing, azimuth timing, orbit knowledge,
attitude knowledge, and localization error introduced by finite resolution,
interpolation, and target detection. These contributions dominate geolocation
bias and localization precision.

The processor branch includes focusing approximations, quantization,
resampling, burst stitching, annotation approximations, and any correction that
is common to the whole product or to a swath or burst. Such terms are often
correlated in space and must therefore be tracked as structured uncertainty
components. In the present product-oriented framing, these processor terms
contribute to three derived observable families:

\begin{itemize}
  \item \textbf{radiometric observables} derived from image intensity, such as
  $|s_{\rm c}|^{2}$, $\sigma^{0}$, noise-equivalent sigma-zero, and radiometric
  uniformity indicators;
  \item \textbf{geometric observables} derived from the image position and
  annotation model, such as target localization, geolocation bias, and range or
  azimuth offsets;
  \item \textbf{phase observables} derived from $\arg(s_{\rm c})$ or local
  phase differences, such as phase bias, phase stability, and phase noise.
\end{itemize}

The figure is intentionally analogous to the uncertainty tree
diagrams used elsewhere in TN3: the calibrated sample is the central product,
while the radiometric, geometric, and phase observables are treated as three
derived measurands with partially shared inputs. Atmospheric and ionospheric
propagation terms are shown explicitly where they affect amplitude, timing, or
phase.

% \begin{figure}[H]
%   \centering
%   \resizebox{0.92\textwidth}{!}{%
%   \begin{tikzpicture}[node distance=1.15cm and 1.15cm, every node/.style={font=\scriptsize}]
%     \node[draw,
%           text width=3.0cm,
%           align=center,
%           minimum height=0.9cm] (raw)
%           {$s_{\rm raw}$};

%     \node[draw,
%           right=1.8cm of raw,
%           text width=2.8cm,
%           align=center,
%           minimum height=0.9cm] (crad)
%           {$C_{\rm rad}$};

%     \node[draw,
%           right=1.8cm of crad,
%           text width=3.2cm,
%           align=center,
%           minimum height=0.9cm] (cgeo)
%           {$C_{\rm geo},\phi_{\rm corr}$};

%     \node[draw,
%           below=1.8cm of crad,
%           text width=6.6cm,
%           align=center,
%           minimum height=1.1cm] (slc)
%           {$s_{\rm c}=f(s_{\rm raw},C_{\rm rad},C_{\rm geo},n,\phi_{\rm corr})+\Delta s$};

%     \node[draw,
%           below=2.0cm of slc,
%           text width=5.6cm,
%           align=center,
%           minimum height=1.1cm] (rad)
%           {$z_{\rm rad}=g_{\rm rad}(s_{\rm c},K_{\rm rad},G,\theta)+\Delta z_{\rm rad}$};

%     \node[draw,
%           below=2.7cm of rad,
%           text width=6.4cm,
%           align=center,
%           minimum height=1.1cm] (geo)
%           {$\mathbf{g}=g_{\rm geo}(p_{\rm img},t_{\rm rg},t_{\rm az},\mathbf{p}_{\rm orb},\mathbf{p}_{\rm att})+\Delta\mathbf{g}$};

%     \node[draw,
%           below=2.7cm of geo,
%           text width=5.6cm,
%           align=center,
%           minimum height=1.1cm] (pha)
%           {$\phi_{\rm L1}=g_{\rm ph}(\arg(s_{\rm c}),\phi_{\rm corr})+\Delta\phi$};

%     \node[draw,
%           left=3.3cm of rad,
%           text width=1.9cm,
%           align=center,
%           minimum height=0.85cm] (dradsc)
%           {$\frac{\partial z_{\rm rad}}{\partial s_{\rm c}}$};

%     \node[draw,
%           right=3.3cm of rad,
%           text width=2.1cm,
%           align=center,
%           minimum height=0.85cm] (dradcrad)
%           {$\frac{\partial z_{\rm rad}}{\partial C_{\rm rad}}$};

%     \node[draw,
%           below left=1.2cm and 2.6cm of rad,
%           text width=1.8cm,
%           align=center,
%           minimum height=0.85cm] (dradn)
%           {$\frac{\partial z_{\rm rad}}{\partial n}$};

%     \node[left=1.6cm of dradsc,
%           text width=2.4cm,
%           align=center] (uradsc)
%           {$u(s_{\rm c})$\\focused\\sample};

%     \node[right=1.6cm of dradcrad,
%           text width=2.4cm,
%           align=center] (ucradleaf)
%           {$u(C_{\rm rad})$\\gain drift};

%     \node[left=1.5cm of dradn,
%           text width=2.4cm,
%           align=center] (unleaf)
%           {$u(n)$\\thermal noise};

%     \node[draw,
%           left=3.4cm of geo,
%           text width=2.0cm,
%           align=center,
%           minimum height=0.85cm] (dgeosc)
%           {$\frac{\partial \mathbf{g}}{\partial s_{\rm c}}$};

%     \node[draw,
%           right=3.4cm of geo,
%           text width=1.9cm,
%           align=center,
%           minimum height=0.85cm] (dgeotrg)
%           {$\frac{\partial \mathbf{g}}{\partial t_{\rm rg}}$};

%     \node[draw,
%           below right=1.2cm and 2.7cm of geo,
%           text width=2.3cm,
%           align=center,
%           minimum height=0.85cm] (dgeoorb)
%           {$\frac{\partial \mathbf{g}}{\partial \mathbf{p}_{\rm orb}}$};

%     \node[left=1.5cm of dgeosc,
%           text width=2.4cm,
%           align=center] (ugeosc)
%           {$u(s_{\rm c})$\\target\\position};

%     \node[right=1.5cm of dgeotrg,
%           text width=2.1cm,
%           align=center] (utrgleaf)
%           {$u(t_{\rm rg})$\\timing};

%     \node[right=1.4cm of dgeoorb,
%           text width=2.4cm,
%           align=center] (uorbleaf)
%           {$u(\mathbf{p}_{\rm orb})$\\orbit state};

%     \node[draw,
%           left=3.3cm of pha,
%           text width=1.9cm,
%           align=center,
%           minimum height=0.85cm] (dphisc)
%           {$\frac{\partial \phi_{\rm L1}}{\partial s_{\rm c}}$};

%     \node[draw,
%           right=3.3cm of pha,
%           text width=2.2cm,
%           align=center,
%           minimum height=0.85cm] (dphicorr)
%           {$\frac{\partial \phi_{\rm L1}}{\partial \phi_{\rm corr}}$};

%     \node[draw,
%           below left=1.2cm and 2.6cm of pha,
%           text width=1.8cm,
%           align=center,
%           minimum height=0.85cm] (dphan)
%           {$\frac{\partial \phi_{\rm L1}}{\partial n}$};

%     \node[left=1.6cm of dphisc,
%           text width=2.4cm,
%           align=center] (uphasc)
%           {$u(s_{\rm c})$\\phase\\sample};

%     \node[right=1.5cm of dphicorr,
%           text width=2.5cm,
%           align=center] (uphicorrleaf)
%           {$u(\phi_{\rm corr})$\\phase correction};

%     \node[left=1.5cm of dphan,
%           text width=2.4cm,
%           align=center] (uphanleaf)
%           {$u(n)$\\noise floor};

%     \draw[-latex] (raw) -- (slc);
%     \draw[-latex] (crad) -- (slc);
%     \draw[-latex] (cgeo) -- (slc);

%     \draw[-latex] (slc) -- (rad);
%     \draw[-latex] (slc) -- (geo);
%     \draw[-latex] (slc) -- (pha);

%     \draw[-latex] (dradsc) -- (rad);
%     \draw[-latex] (dradcrad) -- (rad);
%     \draw[-latex] (dradn) -- (rad);
%     \draw[-latex] (uradsc) -- (dradsc);
%     \draw[-latex] (ucradleaf) -- (dradcrad);
%     \draw[-latex] (unleaf) -- (dradn);

%     \draw[-latex] (dgeosc) -- (geo);
%     \draw[-latex] (dgeotrg) -- (geo);
%     \draw[-latex] (dgeoorb) -- (geo);
%     \draw[-latex] (ugeosc) -- (dgeosc);
%     \draw[-latex] (utrgleaf) -- (dgeotrg);
%     \draw[-latex] (uorbleaf) -- (dgeoorb);

%     \draw[-latex] (dphisc) -- (pha);
%     \draw[-latex] (dphicorr) -- (pha);
%     \draw[-latex] (dphan) -- (pha);
%     \draw[-latex] (uphasc) -- (dphisc);
%     \draw[-latex] (uphicorrleaf) -- (dphicorr);
%     \draw[-latex] (uphanleaf) -- (dphan);
%   \end{tikzpicture}%
%   }
%   \caption{Traceability diagram for a calibrated single-look complex product and the three main families of derived Level-1 observables: radiometric, geometric, and phase observables.}
%   \label{fig:slc_traceability}
% \end{figure}


\begin{figure}[H]
  \centering
  \resizebox{0.92\textwidth}{!}{%
  \begin{tikzpicture}[node distance=1.15cm and 1.15cm, every node/.style={font=\scriptsize}]
    \node[draw,
          text width=3.0cm,
          align=center,
          minimum height=0.9cm] (raw)
          {$s_{\rm raw}$};

    \node[draw,
          right=1.8cm of raw,
          text width=2.8cm,
          align=center,
          minimum height=0.9cm] (crad)
          {$C_{\rm rad}$};

    \node[draw,
          right=1.8cm of crad,
          text width=3.2cm,
          align=center,
          minimum height=0.9cm] (cgeo)
          {$C_{\rm geo},\phi_{\rm corr}$};

    \node[draw,
          above=1.1cm of cgeo,
          text width=3.4cm,
          align=center,
          minimum height=0.9cm] (prop)
          {$\Delta_{\rm trop},\Delta_{\rm iono}$\\path delay, attenuation, Faraday};

    \node[draw,
          below=1.8cm of crad,
          text width=6.6cm,
          align=center,
          minimum height=1.1cm] (slc)
          {$s_{\rm c}=f(s_{\rm raw},C_{\rm rad},C_{\rm geo},n,\phi_{\rm corr})+\Delta s$};

    \node[draw,
          below=2.0cm of slc,
          text width=5.6cm,
          align=center,
          minimum height=1.1cm] (rad)
          {$z_{\rm rad}=g_{\rm rad}(s_{\rm c},K_{\rm rad},G,\theta)+\Delta z_{\rm rad}$};

    \node[draw,
          below=2.7cm of rad,
          text width=6.4cm,
          align=center,
          minimum height=1.1cm] (geo)
          {$\mathbf{g}=g_{\rm geo}(p_{\rm img},t_{\rm rg},t_{\rm az},\mathbf{p}_{\rm orb},\mathbf{p}_{\rm att})+\Delta\mathbf{g}$};

    \node[draw,
          below=2.7cm of geo,
          text width=5.6cm,
          align=center,
          minimum height=1.1cm] (pha)
          {$\phi_{\rm L1}=g_{\rm ph}(\arg(s_{\rm c}),\phi_{\rm corr})+\Delta\phi$};

    \node[draw,
          left=3.3cm of rad,
          text width=1.9cm,
          align=center,
          minimum height=0.85cm] (dradsc)
          {$\frac{\partial z_{\rm rad}}{\partial s_{\rm c}}$};

    \node[draw,
          right=3.3cm of rad,
          text width=2.1cm,
          align=center,
          minimum height=0.85cm] (dradcrad)
          {$\frac{\partial z_{\rm rad}}{\partial C_{\rm rad}}$};

    \node[draw,
          below left=1.2cm and 2.6cm of rad,
          text width=1.8cm,
          align=center,
          minimum height=0.85cm] (dradn)
          {$\frac{\partial z_{\rm rad}}{\partial n}$};

    \node[left=1.6cm of dradsc,
          text width=2.4cm,
          align=center] (uradsc)
          {$u(s_{\rm c})$\\focused\\sample};

    \node[right=1.6cm of dradcrad,
          text width=2.4cm,
          align=center] (ucradleaf)
          {$u(C_{\rm rad})$\\gain drift};

    \node[left=1.5cm of dradn,
          text width=2.4cm,
          align=center] (unleaf)
          {$u(n)$\\thermal noise};

    \node[draw,
          left=3.4cm of geo,
          text width=2.0cm,
          align=center,
          minimum height=0.85cm] (dgeosc)
          {$\frac{\partial \mathbf{g}}{\partial s_{\rm c}}$};

    \node[draw,
          right=3.4cm of geo,
          text width=1.9cm,
          align=center,
          minimum height=0.85cm] (dgeotrg)
          {$\frac{\partial \mathbf{g}}{\partial t_{\rm rg}}$};

    \node[draw,
          below right=1.2cm and 2.7cm of geo,
          text width=2.3cm,
          align=center,
          minimum height=0.85cm] (dgeoorb)
          {$\frac{\partial \mathbf{g}}{\partial \mathbf{p}_{\rm orb}}$};

    \node[draw,
          above right=1.2cm and 2.7cm of geo,
          text width=2.3cm,
          align=center,
          minimum height=0.85cm] (dgeoprop)
          {$\frac{\partial \mathbf{g}}{\partial \Delta_{\rm prop}}$};

    \node[left=1.5cm of dgeosc,
          text width=2.4cm,
          align=center] (ugeosc)
          {$u(s_{\rm c})$\\target\\position};

    \node[right=1.5cm of dgeotrg,
          text width=2.1cm,
          align=center] (utrgleaf)
          {$u(t_{\rm rg})$\\timing};

    \node[right=1.4cm of dgeoorb,
          text width=2.4cm,
          align=center] (uorbleaf)
          {$u(\mathbf{p}_{\rm orb})$\\orbit state};

    \node[right=1.4cm of dgeoprop,
          text width=2.7cm,
          align=center] (upropgeo)
          {$u(\Delta_{\rm prop})$\\tropo/iono delay};

    \node[draw,
          left=3.3cm of pha,
          text width=1.9cm,
          align=center,
          minimum height=0.85cm] (dphisc)
          {$\frac{\partial \phi_{\rm L1}}{\partial s_{\rm c}}$};

    \node[draw,
          right=3.3cm of pha,
          text width=2.2cm,
          align=center,
          minimum height=0.85cm] (dphicorr)
          {$\frac{\partial \phi_{\rm L1}}{\partial \phi_{\rm corr}}$};

    \node[draw,
          below left=1.2cm and 2.6cm of pha,
          text width=1.8cm,
          align=center,
          minimum height=0.85cm] (dphan)
          {$\frac{\partial \phi_{\rm L1}}{\partial n}$};

    \node[draw,
          below right=1.2cm and 2.6cm of pha,
          text width=2.1cm,
          align=center,
          minimum height=0.85cm] (dphiprop)
          {$\frac{\partial \phi_{\rm L1}}{\partial \Delta_{\rm iono}}$};

    \node[left=1.6cm of dphisc,
          text width=2.4cm,
          align=center] (uphasc)
          {$u(s_{\rm c})$\\phase\\sample};

    \node[right=1.5cm of dphicorr,
          text width=2.5cm,
          align=center] (uphicorrleaf)
          {$u(\phi_{\rm corr})$\\phase correction};

    \node[left=1.5cm of dphan,
          text width=2.4cm,
          align=center] (uphanleaf)
          {$u(n)$\\noise floor};

    \node[right=1.5cm of dphiprop,
          text width=2.6cm,
          align=center] (uphiprop)
          {ionospheric phase\\Faraday rotation};

    \draw[-latex] (raw) -- (slc);
    \draw[-latex] (crad) -- (slc);
    \draw[-latex] (cgeo) -- (slc);
    \draw[-latex] (prop) -- (cgeo);

    \draw[-latex] (slc) -- (rad);
    \draw[-latex] (slc) -- (geo);
    \draw[-latex] (slc) -- (pha);

    \draw[-latex] (dradsc) -- (rad);
    \draw[-latex] (dradcrad) -- (rad);
    \draw[-latex] (dradn) -- (rad);
    \draw[-latex] (uradsc) -- (dradsc);
    \draw[-latex] (ucradleaf) -- (dradcrad);
    \draw[-latex] (unleaf) -- (dradn);

    \draw[-latex] (dgeosc) -- (geo);
    \draw[-latex] (dgeotrg) -- (geo);
    \draw[-latex] (dgeoorb) -- (geo);
    \draw[-latex] (dgeoprop) -- (geo);
    \draw[-latex] (ugeosc) -- (dgeosc);
    \draw[-latex] (utrgleaf) -- (dgeotrg);
    \draw[-latex] (uorbleaf) -- (dgeoorb);
    \draw[-latex] (upropgeo) -- (dgeoprop);

    \draw[-latex] (dphisc) -- (pha);
    \draw[-latex] (dphicorr) -- (pha);
    \draw[-latex] (dphan) -- (pha);
    \draw[-latex] (dphiprop) -- (pha);
    \draw[-latex] (uphasc) -- (dphisc);
    \draw[-latex] (uphicorrleaf) -- (dphicorr);
    \draw[-latex] (uphanleaf) -- (dphan);
    \draw[-latex] (uphiprop) -- (dphiprop);
  \end{tikzpicture}%
  }
  \caption{Traceability diagram for a calibrated single-look complex product and the three main families of derived Level-1 observables: radiometric, geometric, and phase observables.}
  \label{fig:slc_traceability}
\end{figure}

\subsubsection{Step 3: Propagation to Level-1 observables}

The complex sample itself is usually not the only reported quantity. Several
Level-1 observables are derived from it, and each inherits a different subset
of the same uncertainty budget. Examples are:
\begin{itemize}
  \item intensity or backscatter-related quantities, derived from $|s_{\rm c}|^2$;
  \item point-target localization coordinates, derived from the image position
  and annotation model;
  \item phase-based observables, derived from $\arg(s_{\rm c})$ or phase
  differences.
\end{itemize}

Accordingly, the uncertainty propagation should not stop at the complex sample,
but continue to the actual Level-1 quantity being reported. For a backscatter
observable $z = |s_{\rm c}|^2$, the uncertainty is driven mainly by the
radiometric branch; for a localization observable $\mathbf{g}=h(s_{\rm c})$,
the geometric branch is dominant. The same single-look complex (SLC) product thus supports
multiple product-level uncertainty statements. A convenient formalization is
to treat the three observable families separately:
\begin{align}
  u^{2}(z_{\rm rad}) &\approx \mathbf{J}_{\rm rad}\, \mathbf{U}_{\rm rad}\,
  \mathbf{J}_{\rm rad}^{\mathsf{T}}, \\
  \mathbf{U}_{\mathbf{g}} &\approx \mathbf{J}_{\rm geo}\, \mathbf{U}_{\rm geo}\,
  \mathbf{J}_{\rm geo}^{\mathsf{T}}, \\
  u^{2}(\phi_{\rm L1}) &\approx \mathbf{J}_{\rm ph}\, \mathbf{U}_{\rm ph}\,
  \mathbf{J}_{\rm ph}^{\mathsf{T}},
  \label{eq:slc_three_branches}
\end{align}
where each covariance matrix contains the dominant inputs for the corresponding
observable family. This makes it possible to report a product-consistent yet
observable-specific uncertainty estimate.

\subsubsection{Step 4: Uncertainty effects table for the single-look complex product}

Table~\ref{tab:slc_effects_table} summarizes the main effects contributing to
the uncertainty of the three observable families derived from the single-look
complex product. In contrast with the Doppler example, the purpose here is not
to quantify one single scalar measurand, but to identify which sources are
dominant for radiometric, geometric, and phase-related uses of the same
product.

\begin{table*}[ht]
  \scriptsize
  \centering
  \caption{Uncertainty effects table for the calibrated single-look complex product and its main derived Level-1 observables.}
  \resizebox{\textwidth}{!}{%
  \begin{tabular}{|p{2.8cm}|c|c|c|c|c|c|p{2.4cm}|}
    \hline
    \textbf{Name of effect} & \textbf{Radiometric} & \textbf{Geometric} & \textbf{Phase} & \textbf{Maturity of} & \textbf{Correlation} & \textbf{Sensitivity} & \textbf{Typical uncertainty form} \\ 
     & \textbf{observable} & \textbf{observable} & \textbf{observable} & \textbf{estimate} & \textbf{type/scale} & \textbf{coefficient} &  \\ \hline
    Absolute calibration factor & Yes & No & No & 2--3 & systematic / swath & high for $z_{\rm rad}$ & multiplicative bias \\ \hline
    Thermal-noise subtraction & Yes & No & Yes & 2 & random + structured / range & medium & additive noise term \\ \hline
    Antenna pattern normalization & Yes & No & No & 2--3 & structured / range & high & range-dependent bias \\ \hline
    Speckle and finite-sample statistics & Yes & weak & weak & 3 & random / local window & high locally & variance reduction with averaging \\ \hline
    Range and azimuth timing & No & Yes & weak & 2--3 & systematic / product & high for $\mathbf{g}$ & additive localization bias \\ \hline
    Orbit and attitude knowledge & No & Yes & Yes & 2--3 & structured / scene-product & high for $\mathbf{g}$ & covariance term \\ \hline
    Target localization/interpolation & No & Yes & No & 2 & random / local target & medium--high & estimator variance \\ \hline
    Oscillator and phase correction terms & No & No & Yes & 1--2 & structured / burst-product & high for $\phi_{\rm L1}$ & additive phase bias/noise \\ \hline
    Focusing and resampling approximations & Yes & Yes & Yes & 1--2 & structured / burst-swath & medium & processor-dependent residual \\ \hline
    Quantization and digitization & weak & weak & weak & 1--2 & random / sample & low--medium & additive noise \\ \hline
  \end{tabular}
  }
  \label{tab:slc_effects_table}
\end{table*}

\subsubsection{Step 5: Product-level reporting}

For a calibrated single-look complex (SLC), it is therefore recommended to report uncertainty in
terms of product-level quality indicators, for example:
\begin{itemize}
  \item radiometric accuracy and stability for backscatter-related quantities;
  \item localization accuracy and bias for geometric quality;
  \item phase stability for phase-sensitive applications.
\end{itemize}
The key point is that these are all uncertainty statements of the \emph{same}
Level-1 product, not separate products created by the calibration procedures.
In practice, this means that a single-look complex product should ideally be
delivered with either dedicated quality layers or documentation that separates:
\begin{itemize}
  \item uncertainties attached to radiometric observables derived from the
  product;
  \item uncertainties attached to geometric observables derived from the
  product; and
  \item uncertainties attached to phase observables derived from the product.
\end{itemize}
This separation mirrors the way users actually exploit Level-1 data and makes
the uncertainty statement much more actionable for downstream applications.

\subsection{Polarimetric Level-1 product}

For a polarimetric Level-1 product, the measurand is no longer a single complex
sample but the calibrated scattering matrix, or equivalently the covariance or
coherency representation derived from it. All single-channel SLC uncertainty effects still
apply to every polarimetric channel. The additional polarimetric information
allows some channel imbalance, cross-talk, co-polar phase, and Faraday-rotation
effects to be estimated or compensated, but it also introduces covariance
between channels that must be retained in the uncertainty budget.

\subsubsection{Step 1: Definition of measurand and measurement model}

Let $\mathbf{S}_{\rm c}$ denote the calibrated scattering matrix. A simplified
model is
\begin{equation}
  \mathbf{S}_{\rm c} = \mathbf{M}_{\rm r}^{-1}
  \mathbf{S}_{\rm m}\, \mathbf{M}_{\rm t}^{-1} + \Delta\mathbf{S},
  \label{eq:l1_pol_product_model}
\end{equation}
where $\mathbf{S}_{\rm m}$ is the measured matrix,
$\mathbf{M}_{\rm r}$ and $\mathbf{M}_{\rm t}$ summarize receive and transmit
channel distortions, and $\Delta\mathbf{S}$ is the residual polarimetric error.





\subsubsection{Step 2: Traceability diagram for the polarimetric product}

For the polarimetric product, the uncertainty contributions arise from:
\begin{itemize}
  \item the radiometric branch inherited from each polarization channel;
  \item the polarimetric calibration branch, including channel imbalance,
  cross-talk, phase offsets, and possible Faraday-rotation-related effects;
  \item the geometric co-registration branch between channels;
  \item the statistical branch, associated with covariance estimation over a
  finite number of samples.
\end{itemize}

Figure~\ref{fig:pol_traceability} shows a compact uncertainty tree for the
polarimetric Level-1 product. The calibrated scattering matrix is the central
measurand, while channel imbalance, cross-talk, and co-polar phase observables
are treated as derived quantities supported by reference targets and channel
statistics.

% \begin{figure}[H]
%   \centering
%   \resizebox{0.96\textwidth}{!}{%
%   \begin{tikzpicture}[node distance=1.25cm and 1.2cm, every node/.style={font=\small}]
%     \node[draw,text width=2.8cm,align=center,minimum height=0.9cm] (sm)
%       {$\mathbf{S}_{\rm m}$};
%     \node[draw,right=1.9cm of sm,text width=2.3cm,align=center,minimum height=0.9cm] (mr)
%       {$\mathbf{M}_{\rm r}$};
%     \node[draw,right=1.9cm of mr,text width=2.3cm,align=center,minimum height=0.9cm] (mt)
%       {$\mathbf{M}_{\rm t}$};

%     \node[draw,below=1.6cm of mr,text width=6.8cm,align=center,minimum height=1.05cm] (sc)
%       {$\mathbf{S}_{\rm c}=\mathbf{M}_{\rm r}^{-1}\mathbf{S}_{\rm m}\mathbf{M}_{\rm t}^{-1}+\Delta\mathbf{S}$};

%     \node[draw,below=2.0cm of sc,text width=4.8cm,align=center,minimum height=1.0cm] (imb)
%       {$q_{\rm imb}=g_{\rm imb}(\mathbf{S}_{\rm c})+\Delta q_{\rm imb}$};
%     \node[draw,below=2.4cm of imb,text width=4.8cm,align=center,minimum height=1.0cm] (xt)
%       {$q_{\rm xt}=g_{\rm xt}(\mathbf{S}_{\rm c})+\Delta q_{\rm xt}$};
%     \node[draw,below=2.4cm of xt,text width=4.8cm,align=center,minimum height=1.0cm] (phb)
%       {$q_{\phi}=g_{\phi}(\mathbf{S}_{\rm c})+\Delta q_{\phi}$};

%     \node[draw,left=3.0cm of imb,text width=1.9cm,align=center,minimum height=0.85cm] (dimb)
%       {$\frac{\partial q_{\rm imb}}{\partial \mathbf{S}_{\rm c}}$};
%     \node[left=1.3cm of dimb,text width=2.2cm,align=center] (uimb)
%       {$u(\mathbf{M}_{\rm r,t})$\\channel imbalance};

%     \node[draw,left=3.0cm of xt,text width=1.9cm,align=center,minimum height=0.85cm] (dxt)
%       {$\frac{\partial q_{\rm xt}}{\partial \mathbf{S}_{\rm c}}$};
%     \node[left=1.3cm of dxt,text width=2.0cm,align=center] (uxt)
%       {$u(\mathrm{XT})$\\cross-talk};

%     \node[draw,right=3.0cm of phb,text width=1.9cm,align=center,minimum height=0.85cm] (dphb)
%       {$\frac{\partial q_{\phi}}{\partial \mathbf{S}_{\rm c}}$};
%     \node[right=1.3cm of dphb,text width=2.2cm,align=center] (uphb)
%       {$u(\Delta\phi_{\rm co})$\\co-pol phase offset};

%     \draw[-latex] (sm) -- (sc);
%     \draw[-latex] (mr) -- (sc);
%     \draw[-latex] (mt) -- (sc);
%     \draw[-latex] (sc) -- (imb);
%     \draw[-latex] (sc) -- (xt);
%     \draw[-latex] (sc) -- (phb);
%     \draw[-latex] (dimb) -- (imb);
%     \draw[-latex] (uimb) -- (dimb);
%     \draw[-latex] (dxt) -- (xt);
%     \draw[-latex] (uxt) -- (dxt);
%     \draw[-latex] (dphb) -- (phb);
%     \draw[-latex] (uphb) -- (dphb);
%   \end{tikzpicture}%
%   }
%   \caption{Traceability diagram for the polarimetric Level-1 product and the main observables related to channel imbalance, cross-talk, and co-polar phase bias.}
%   \label{fig:pol_traceability}
% \end{figure}
\begin{figure}[H]
  \centering
  \resizebox{0.96\textwidth}{!}{%
  \begin{tikzpicture}[node distance=1.25cm and 1.2cm, every node/.style={font=\small}]
    \node[draw,text width=2.8cm,align=center,minimum height=0.9cm] (sm)
      {$\mathbf{S}_{\rm m}$};
    \node[draw,right=1.9cm of sm,text width=2.3cm,align=center,minimum height=0.9cm] (mr)
      {$\mathbf{M}_{\rm r}$};
    \node[draw,right=1.9cm of mr,text width=2.3cm,align=center,minimum height=0.9cm] (mt)
      {$\mathbf{M}_{\rm t}$};

    \node[draw,below=1.6cm of mr,text width=6.8cm,align=center,minimum height=1.05cm] (sc)
      {$\mathbf{S}_{\rm c}=\mathbf{M}_{\rm r}^{-1}\mathbf{S}_{\rm m}\mathbf{M}_{\rm t}^{-1}+\Delta\mathbf{S}$};

    \node[draw,below=2.0cm of sc,text width=4.8cm,align=center,minimum height=1.0cm] (imb)
      {$q_{\rm imb}=g_{\rm imb}(\mathbf{S}_{\rm c})+\Delta q_{\rm imb}$};
    \node[draw,below=2.4cm of imb,text width=4.8cm,align=center,minimum height=1.0cm] (xt)
      {$q_{\rm xt}=g_{\rm xt}(\mathbf{S}_{\rm c})+\Delta q_{\rm xt}$};
    \node[draw,below=2.4cm of xt,text width=4.8cm,align=center,minimum height=1.0cm] (phb)
      {$q_{\phi}=g_{\phi}(\mathbf{S}_{\rm c})+\Delta q_{\phi}$};

    \node[draw,left=3.0cm of imb,text width=1.9cm,align=center,minimum height=0.85cm] (dimb)
      {$\frac{\partial q_{\rm imb}}{\partial \mathbf{S}_{\rm c}}$};
    \node[left=1.3cm of dimb,text width=2.2cm,align=center] (uimb)
      {$u(\mathbf{M}_{\rm r,t})$\\channel imbalance};

    \node[draw,left=3.0cm of xt,text width=1.9cm,align=center,minimum height=0.85cm] (dxt)
      {$\frac{\partial q_{\rm xt}}{\partial \mathbf{S}_{\rm c}}$};
    \node[left=1.3cm of dxt,text width=2.0cm,align=center] (uxt)
      {$u(\mathrm{XT})$\\cross-talk};

    \node[draw,right=3.0cm of phb,text width=1.9cm,align=center,minimum height=0.85cm] (dphb)
      {$\frac{\partial q_{\phi}}{\partial \mathbf{S}_{\rm c}}$};
    \node[right=1.3cm of dphb,text width=2.2cm,align=center] (uphb)
      {$u(\Delta\phi_{\rm co})$\\co-pol phase offset};

    \node[draw,left=3.0cm of phb,text width=2.0cm,align=center,minimum height=0.85cm] (dpoliono)
      {$\frac{\partial q_{\phi}}{\partial F_{\rm iono}}$};
    \node[left=1.3cm of dpoliono,text width=2.4cm,align=center] (uiono)
      {$u(F_{\rm iono})$\\Faraday rotation};

    \draw[-latex] (sm) -- (sc);
    \draw[-latex] (mr) -- (sc);
    \draw[-latex] (mt) -- (sc);
    \draw[-latex] (sc) -- (imb);
    \draw[-latex] (sc) -- (xt);
    \draw[-latex] (sc) -- (phb);
    \draw[-latex] (dimb) -- (imb);
    \draw[-latex] (uimb) -- (dimb);
    \draw[-latex] (dxt) -- (xt);
    \draw[-latex] (uxt) -- (dxt);
    \draw[-latex] (dphb) -- (phb);
    \draw[-latex] (uphb) -- (dphb);
    \draw[-latex] (dpoliono) -- (phb);
    \draw[-latex] (uiono) -- (dpoliono);
  \end{tikzpicture}%
  }
  \caption{Traceability diagram for the polarimetric Level-1 product and the main observables related to channel imbalance, cross-talk, and co-polar phase bias.}
  \label{fig:pol_traceability}
\end{figure}

\subsubsection{Step 3: Uncertainty effects table for the polarimetric product}
This is provided in Table \ref{tab:pol_effects_table}.

\begin{table*}[ht]
  \scriptsize
  \centering
  \caption{Uncertainty effects table for the polarimetric Level-1 product and its main observables.}
  \resizebox{\textwidth}{!}{%
  \begin{tabular}{|p{3.0cm}|c|c|c|c|c|p{2.8cm}|}
    \hline
    \textbf{Name of effect} & \textbf{Imbalance} & \textbf{Cross-talk} & \textbf{Phase bias} & \textbf{Maturity} & \textbf{Correlation} & \textbf{Typical uncertainty form} \\ \hline
    Channel gain imbalance & Yes & weak & weak & 2--3 & structured / channel & multiplicative bias \\ \hline
    Cross-talk coefficients & weak & Yes & weak & 2--3 & structured / channel & leakage term \\ \hline
    Co-polar phase offset & No & weak & Yes & 2--3 & structured / product & additive phase bias \\ \hline
    Reference target response & Yes & Yes & Yes & 2 & scene-dependent / target & calibration reference uncertainty \\ \hline
    Channel co-registration & weak & Yes & Yes & 2 & structured / image & residual misregistration \\ \hline
    Thermal noise/statistics & weak & Yes & Yes & 2--3 & random / window & covariance estimation variance \\ \hline
    Faraday/ionospheric effects & weak & weak & Yes & 1--2 & scene-dependent & propagation-induced phase term \\ \hline
  \end{tabular}
  }
  \label{tab:pol_effects_table}
\end{table*}

\subsubsection{Step 4: Propagation to polarimetric observables}

Product-level observables include co-polar ratio, cross-talk indicators,
co-polar phase difference, and decomposition variables derived from covariance
or coherency matrices. For any such observable $q = g(\mathbf{S}_{\rm c})$,
the uncertainty propagation follows the covariance of the real and imaginary
parts of the calibrated matrix elements. Correlations between channels are
usually important and should be retained in the analysis.

\subsubsection{Step 5: Product-level reporting}

For a polarimetric single-look complex (SLC), the reported uncertainty should be associated with
product observables, such as residual channel imbalance, residual cross-talk,
co-polar phase bias, or reciprocity diagnostics, rather than with the
calibration procedure in isolation.

\subsection{Interferometric Level-1 product}

An interferometric product is built from two calibrated and co-registered
complex Level-1 products. Even when the product classification differs across
missions, the measurands of interest are typically the interferometric phase,
coherence, and related quality indicators.

\subsubsection{Step 1: Definition of measurand and measurement model}

For two calibrated single-look complex (SLC) products $s_{1}$ and $s_{2}$, the basic
interferometric observables are
\begin{equation}
  I = s_{1}s_{2}^{\ast},
  \qquad
  \hat{\phi} = \arg(I),
  \qquad
  \hat{\gamma} =
  \frac{\left|\sum_{k=1}^{N} s_{1,k}s_{2,k}^{\ast}\right|}
  {\sqrt{\sum_{k=1}^{N}|s_{1,k}|^2\sum_{k=1}^{N}|s_{2,k}|^2}}.
  \label{eq:l1_insar_model}
\end{equation}
The measurement model therefore contains both input-product uncertainties and
interferometric-estimator uncertainties.



\subsubsection{Step 2: Traceability diagram for the interferometric product}

The uncertainty budget of the interferometric product contains:
\begin{itemize}
  \item uncertainties inherited from each input single-look complex (SLC);
  \item co-registration uncertainty in range and azimuth;
  \item baseline, orbit, and attitude uncertainty;
  \item thermal-noise and finite-window estimator uncertainty;
  \item decorrelation terms of temporal, geometric, or volumetric origin;
  \item optional unwrapping uncertainty if unwrapped phase is analysed.
\end{itemize}

When coherence or phase is estimated after flattening or topographic
demodulation, the demodulation phase is itself an input quantity. The
uncertainty of the external DEM, geoid or height reference, baseline, and orbit
therefore propagates into the residual phase and, in steep or heterogeneous
terrain, can also affect coherence through imperfect phase removal within the
estimation window. For many low-relief coherence applications this DEM-driven
term may be negligible, but it shall be documented as an auxiliary-data
contribution rather than omitted by assumption.


Figure~\ref{fig:insar_traceability} shows the uncertainty tree for the
interferometric product. The two calibrated input single-look complex products
feed the interferogram, from which phase and coherence are estimated as the
main Level-1 observables.

\begin{figure}[H]
  \centering
  \resizebox{0.96\textwidth}{!}{%
  \begin{tikzpicture}[node distance=1.25cm and 1.2cm, every node/.style={font=\small}]
    \node[draw,text width=2.6cm,align=center,minimum height=0.9cm] (s1)
      {$s_{1}$};
    \node[draw,right=2.0cm of s1,text width=2.6cm,align=center,minimum height=0.9cm] (s2)
      {$s_{2}$};

    \node[draw,below=1.7cm of $(s1)!0.5!(s2)$,text width=5.8cm,align=center,minimum height=1.0cm] (igram)
      {$I=s_{1}s_{2}^{\ast}+\Delta I$};

    \node[draw,below=2.0cm of igram,text width=4.6cm,align=center,minimum height=1.0cm] (coh)
      {$\hat{\gamma}=g_{\gamma}(I,N)+\Delta\gamma$};
    \node[draw,below=2.3cm of coh,text width=4.6cm,align=center,minimum height=1.0cm] (iph)
      {$\hat{\phi}=g_{\phi}(I,\mathbf{b})+\Delta\phi$};

    \node[draw,left=2.8cm of coh,text width=1.8cm,align=center,minimum height=0.85cm] (dcoh)
      {$\frac{\partial \hat{\gamma}}{\partial I}$};
    \node[left=1.3cm of dcoh,text width=2.5cm,align=center] (ucoh)
      {$u(N)$\\window/noise};

    \node[draw,right=2.8cm of iph,text width=1.8cm,align=center,minimum height=0.85cm] (diph)
      {$\frac{\partial \hat{\phi}}{\partial I}$};
    \node[right=1.3cm of diph,text width=2.6cm,align=center] (uiph)
      {$u(\mathbf{b})$\\baseline/orbit};

    \node[draw,left=2.8cm of iph,text width=2.0cm,align=center,minimum height=0.85cm] (dinsprop)
      {$\frac{\partial \hat{\phi}}{\partial \Delta L_{\rm prop}}$};
    \node[left=1.3cm of dinsprop,text width=2.7cm,align=center] (uinsprop)
      {tropospheric and\\ionospheric delay};

    \node[draw,left=2.8cm of igram,text width=1.8cm,align=center,minimum height=0.85cm] (dig)
      {$\frac{\partial I}{\partial s_{1,2}}$};
    \node[left=1.3cm of dig,text width=2.3cm,align=center] (uslc)
      {$u(s_{1}),u(s_{2})$\\input SLCs};

    \draw[-latex] (s1) -- (igram);
    \draw[-latex] (s2) -- (igram);
    \draw[-latex] (dig) -- (igram);
    \draw[-latex] (uslc) -- (dig);
    \draw[-latex] (igram) -- (coh);
    \draw[-latex] (igram) -- (iph);
    \draw[-latex] (dcoh) -- (coh);
    \draw[-latex] (ucoh) -- (dcoh);
    \draw[-latex] (diph) -- (iph);
    \draw[-latex] (uiph) -- (diph);
    \draw[-latex] (dinsprop) -- (iph);
    \draw[-latex] (uinsprop) -- (dinsprop);
  \end{tikzpicture}%
  }
  \caption{Traceability diagram for the interferometric Level-1 product and its principal observables: coherence and interferometric phase.}
  \label{fig:insar_traceability}
\end{figure}

\subsubsection{Step 3: Uncertainty effects table for the interferometric product}
This is reported in Table \ref{tab:insar_effects_table}
\begin{table*}[ht]
  \scriptsize
  \centering
  \caption{Uncertainty effects table for the interferometric Level-1 product and its main observables.}
  \resizebox{\textwidth}{!}{%
  \begin{tabular}{|p{3.0cm}|c|c|c|c|p{3.2cm}|}
    \hline
    \textbf{Name of effect} & \textbf{Coherence} & \textbf{Phase} & \textbf{Maturity} & \textbf{Correlation} & \textbf{Typical uncertainty form} \\ \hline
    Input SLC radiometric/phase errors & Yes & Yes & 2--3 & structured / image pair & inherited covariance term \\ \hline
    Co-registration error & Yes & Yes & 2--3 & structured / burst-scene & residual decorrelation / phase ramp \\ \hline
    Baseline and orbit uncertainty & weak & Yes & 2--3 & structured / pair & geometric phase term \\ \hline
    Thermal noise and finite window & Yes & Yes & 3 & random / estimator window & estimator variance \\ \hline
    Temporal decorrelation & Yes & weak & 2 & scene-dependent & coherence loss \\ \hline
    Geometric/volume decorrelation & Yes & weak & 2 & scene-dependent & coherence loss \\ \hline
    Phase unwrapping (if used) & No & Yes & 1--2 & structured / connected region & discrete ambiguity term \\ \hline
  \end{tabular}
  }
  \label{tab:insar_effects_table}
\end{table*}

\subsubsection{Step 4: Propagation and interpretation}

The uncertainty of coherence and interferometric phase is strongly estimator-
dependent and scene-dependent. For this reason, a mixed strategy is often more
appropriate than a purely analytical one: the deterministic correction terms
can be propagated analytically, while the statistical behaviour of the
interferometric estimator is better characterized through repeated acquisitions,
homogeneous calibration areas, or Monte Carlo simulation of complex signals.

\subsubsection{Step 5: Product-level reporting}

For an interferometric product, the uncertainty statement should always be tied
to the acquisition pair, the estimator window, and the scene class. In this
way, the reported uncertainty remains a property of the product that the user
actually interprets, rather than a generic statement about one element of the
processing chain.



\subsection{Geocoded ARD/NRB Level-1 product}
\label{sec:l1_nrb_uncertainty}

Normalised Radar Backscatter (NRB) is a Level-1 Analysis Ready Data product in which the calibrated SAR measurement is terrain-corrected, radiometrically normalised, resampled, and delivered on a defined map grid. The CEOS NRB Product Family Specification therefore extends the uncertainty problem beyond the source SLC or detected product: the final measurand is not only a calibrated backscatter value, but a terrain-normalised backscatter value assigned to a specific projected pixel with associated metadata, masks, and ancillary-data lineage \cite{NRB_PFS}.

\subsubsection{Step 1: Definition of measurand and measurement model}

For an NRB pixel, a generic measurand can be written as
\begin{equation}
  y_{\rm NRB}(\mathbf{p}) =
  T_{\rm rad}\!\left(z_{\rm src}, \theta_{\rm loc}, \mathbf{d}_{\rm DEM}, \mathbf{m}\right)
  + \Delta y_{\rm NRB},
  \label{eq:nrb_model}
\end{equation}
where $\mathbf{p}$ is the output map-grid pixel, $z_{\rm src}$ is the calibrated source-product backscatter or intensity quantity, $T_{\rm rad}$ is the radiometric terrain-normalisation and geocoding operator, $\theta_{\rm loc}$ is the local incidence-angle information, $\mathbf{d}_{\rm DEM}$ denotes DEM and terrain variables, $\mathbf{m}$ denotes product masks and validity flags, and $\Delta y_{\rm NRB}$ represents residual modelling and implementation effects. The corresponding geometric measurand is the map location of the output pixel,
\begin{equation}
  \mathbf{g}_{\rm NRB} = G\!\left(\mathbf{g}_{\rm src}, \mathbf{x}_{\rm orb}, \mathbf{d}_{\rm DEM}, \mathbf{x}_{\rm CRS}, \mathbf{x}_{\rm resamp}\right) + \Delta \mathbf{g}_{\rm NRB},
  \label{eq:nrb_geometric_model}
\end{equation}
where $G$ is the geocoding operator, $\mathbf{g}_{\rm src}$ is the source-product geometric information, $\mathbf{x}_{\rm orb}$ represents orbit and attitude information, $\mathbf{x}_{\rm CRS}$ defines the coordinate reference system and map-grid convention, and $\mathbf{x}_{\rm resamp}$ represents interpolation and resampling choices.

\subsubsection{Step 2: Traceability diagram for the NRB product}

The NRB uncertainty budget should be organized into inherited and product-generation terms:
\begin{itemize}
  \item \textbf{Inherited source-product terms:} absolute radiometric calibration, antenna-pattern correction, noise subtraction, timing, orbit, attitude, and any source-product processor approximations.
  \item \textbf{Radiometric terrain-normalisation terms:} DEM height uncertainty, slope/aspect uncertainty, local incidence-angle estimation, scattering-area model, terrain-flattening convention, and residual layover or shadow contamination.
  \item \textbf{Geometric and map-grid terms:} coordinate reference system definition, pixel coordinate convention, grid origin, pixel spacing, orthorectification model, DEM-induced geolocation error, and reference-image or GCP uncertainty used for validation.
  \item \textbf{Resampling and aggregation terms:} interpolation kernel, multilooking or averaging, correlated neighbouring pixels, edge effects, and changes in speckle statistics caused by reprojection.
  \item \textbf{Mask and metadata terms:} valid-data mask, layover/shadow mask, water or no-data masks where applicable, metadata machine readability, ancillary-data versioning, and temporal validity of the source and reference data.
\end{itemize}

Figure~\ref{fig:nrb_traceability} shows the uncertainty tree for a geocoded
ARD/NRB Level-1 product. The diagram separates the inherited source-product
branch from the product-generation branches introduced by terrain
normalisation, geocoding, resampling, masking, and ancillary-data lineage. This
is important because an NRB pixel is not only a calibrated backscatter value;
it is a terrain-normalised value assigned to a defined map-grid pixel and
qualified by masks and metadata.

\begin{figure}[H]
  \centering
  \resizebox{\textwidth}{!}{%
  \begin{tikzpicture}[node distance=1.35cm and 1.55cm, every node/.style={font=\small}]
    \node[draw,text width=3.2cm,align=center,minimum height=0.9cm] (src)
      {$z_{\rm src}$\\source radiometry};
    \node[draw,right=1.8cm of src,text width=3.4cm,align=center,minimum height=0.9cm] (dem)
      {$\mathbf{d}_{\rm DEM},\theta_{\rm loc}$\\terrain geometry};
    \node[draw,right=1.8cm of dem,text width=3.2cm,align=center,minimum height=0.9cm] (geo)
      {$\mathbf{x}_{\rm orb},\mathbf{x}_{\rm CRS}$\\geocoding};
    \node[draw,above=1.45cm of geo,text width=3.8cm,align=center,minimum height=0.9cm] (atm)
      {$\mathbf{x}_{\rm atm},\mathbf{x}_{\rm iono}$\\attenuation, delay};

    \node[draw,below=2.0cm of dem,text width=7.8cm,align=center,minimum height=1.0cm] (trad)
      {$T_{\rm rad}(z_{\rm src},\theta_{\rm loc},\mathbf{d}_{\rm DEM},\mathbf{m})+\Delta y_{\rm NRB}$};

    \node[draw,below=2.0cm of trad,text width=7.3cm,align=center,minimum height=1.0cm] (grid)
      {$y_{\rm NRB}(\mathbf{p}),\;\mathbf{g}_{\rm NRB}=G(\mathbf{g}_{\rm src},\mathbf{x}_{\rm orb},\mathbf{d}_{\rm DEM},\mathbf{x}_{\rm CRS},\mathbf{x}_{\rm resamp})+\Delta\mathbf{g}_{\rm NRB}$};

    \node[draw,left=3.0cm of trad,text width=1.7cm,align=center,minimum height=0.85cm] (dsrc)
      {$\frac{\partial y}{\partial z_{\rm src}}$};
    \node[left=1.45cm of dsrc,text width=2.6cm,align=center] (usrc)
      {$u(z_{\rm src})$\\calibration, noise, speckle};

    \node[draw,above=1.55cm of trad,text width=1.9cm,align=center,minimum height=0.85cm] (ddem)
      {$\frac{\partial y}{\partial \mathbf{d}_{\rm DEM}}$};
    \node[above=1.25cm of ddem,text width=3.0cm,align=center] (udem)
      {$u(\mathbf{d}_{\rm DEM}),u(\theta_{\rm loc})$\\height, slope, aspect};

    \node[draw,right=3.0cm of trad,text width=1.8cm,align=center,minimum height=0.85cm] (dgeo)
      {$\frac{\partial \mathbf{g}}{\partial \mathbf{x}_{\rm geo}}$};
    \node[right=1.45cm of dgeo,text width=2.8cm,align=center] (ugeo)
      {$u(\mathbf{x}_{\rm orb}),u(\mathbf{x}_{\rm CRS})$\\orbit, grid, CRS};

    \node[draw,above right=2.05cm and 2.35cm of trad,text width=2.0cm,align=center,minimum height=0.85cm] (datm)
      {$\frac{\partial y,\mathbf{g}}{\partial \mathbf{x}_{\rm prop}}$};
    \node[right=1.35cm of datm,text width=2.8cm,align=center] (uatm)
      {$u(\mathbf{x}_{\rm atm}),u(\mathbf{x}_{\rm iono})$\\meteo, TEC maps};

    \node[draw,left=3.0cm of grid,text width=1.9cm,align=center,minimum height=0.85cm] (dres)
      {$\frac{\partial y}{\partial \mathbf{x}_{\rm resamp}}$};
    \node[left=1.45cm of dres,text width=2.8cm,align=center] (ures)
      {$u(\mathbf{x}_{\rm resamp})$\\kernel, aggregation};

    \node[draw,right=3.0cm of grid,text width=1.8cm,align=center,minimum height=0.85cm] (dmask)
      {$u_{\rm mask}$};
    \node[right=1.45cm of dmask,text width=2.7cm,align=center] (umask)
      {layover, shadow, no-data, validity flags};

    \draw[-latex] (src) -- (trad);
    \draw[-latex] (dem) -- (trad);
    \draw[-latex] (geo) -- (trad);
    \draw[-latex] (atm) -- (trad);
    \draw[-latex] (trad) -- (grid);
    \draw[-latex] (usrc) -- (dsrc);
    \draw[-latex] (dsrc) -- (trad);
    \draw[-latex] (udem) -- (ddem);
    \draw[-latex] (ddem) -- (trad);
    \draw[-latex] (ugeo) -- (dgeo);
    \draw[-latex] (dgeo) -- (trad);
    \draw[-latex] (uatm) -- (datm);
    \draw[-latex] (datm) -- (trad);
    \draw[-latex] (datm) -- (grid);
    \draw[-latex] (ures) -- (dres);
    \draw[-latex] (dres) -- (grid);
    \draw[-latex] (umask) -- (dmask);
    \draw[-latex] (dmask) -- (grid);
  \end{tikzpicture}%
  }
  \caption{Traceability diagram for a geocoded ARD/NRB Level-1 product. The final measurands are the terrain-normalised backscatter value $y_{\rm NRB}(\mathbf{p})$ and its associated map-grid location $\mathbf{g}_{\rm NRB}$.}
  \label{fig:nrb_traceability}
\end{figure}

\subsubsection{Step 4: Uncertainty effects table for the NRB product}

Table~\ref{tab:nrb_effects_table} summarizes the main uncertainty effects for
NRB/ARD products. The table distinguishes radiometric uncertainty, geometric
uncertainty, and mask/metadata validity because these components have different
correlation structures and are used differently by downstream users.

\begin{table*}[ht]
  \scriptsize
  \centering
  \caption{Uncertainty effects table for the geocoded NRB Level-1 product.}
  \resizebox{\textwidth}{!}{%
  \begin{tabular}{|p{3.2cm}|c|c|c|c|p{3.3cm}|p{3.2cm}|}
    \hline
    \textbf{Name of effect} & \textbf{Radiometric} & \textbf{Geometric} & \textbf{Mask / metadata} & \textbf{Maturity} & \textbf{Correlation type / scale} & \textbf{Typical uncertainty form} \\ \hline
    Inherited source-product calibration & Yes & weak & No & 2--3 & structured / swath-product & multiplicative backscatter bias \\ \hline
    Source-product noise, speckle, and finite looks & Yes & weak & No & 2--3 & random to local-window & variance term reduced by averaging \\ \hline
    DEM height, slope, and aspect & Yes & Yes & weak & 2--3 & spatially correlated / terrain tile & local-incidence and geolocation covariance \\ \hline
    Local incidence-angle estimation & Yes & weak & No & 2--3 & structured / terrain facet & terrain-normalisation sensitivity term \\ \hline
    Terrain-flattening and scattering-area model & Yes & weak & No & 1--2 & structured / scene class & model-convention bias \\ \hline
    Orbit, attitude, and timing inherited by geocoding & weak & Yes & No & 2--3 & structured / acquisition & map-location covariance term \\ \hline
    CRS, grid origin, pixel spacing, and coordinate convention & No & Yes & Yes & 2--3 & product-wide systematic & map-grid convention and metadata term \\ \hline
    Resampling, interpolation, and aggregation & Yes & Yes & weak & 1--2 & local-window / kernel footprint & smoothing, covariance, and edge effects \\ \hline
    Layover, shadow, water, no-data, and validity masks & weak & weak & Yes & 1--2 & object/local-region & classification or flagging uncertainty \\ \hline
    Ancillary-data lineage and processor version & Yes & Yes & Yes & 1--2 & product-wide / processing baseline & traceability and reproducibility term \\ \hline
    Reference image or GCP validation data & weak & Yes & weak & 2 & local to scene-wide & validation-reference covariance \\ \hline
  \end{tabular}
  }
  \label{tab:nrb_effects_table}
\end{table*}

\subsubsection{Step 5: Propagation and covariance structure}

For NRB, uncertainty propagation should account for both radiometric and geometric covariance. In compact form,
\begin{equation}
  u^2\!\left(y_{\rm NRB}\right) \approx
  \mathbf{J}_{T}\,\mathbf{U}_{\rm src}\,\mathbf{J}_{T}^{\mathsf{T}}
  + \mathbf{J}_{a}\,\mathbf{U}_{\rm aux}\,\mathbf{J}_{a}^{\mathsf{T}}
  + \mathbf{J}_{p}\,\mathbf{U}_{\rm proc}\,\mathbf{J}_{p}^{\mathsf{T}}
  + u^2_{\rm mask},
  \label{eq:nrb_uncertainty}
\end{equation}
where $\mathbf{U}_{\rm src}$ contains inherited source-product covariance, $\mathbf{U}_{\rm aux}$ contains DEM, orbit, incidence-angle, and other ancillary-data covariance, $\mathbf{U}_{\rm proc}$ contains geocoding and resampling implementation terms, and $u_{\rm mask}$ represents uncertainty associated with validity classification. These terms are often spatially correlated: for example, a DEM bias or orbit error can affect many neighbouring pixels coherently, whereas speckle and feature-identification errors may be more local. A GUM-compliant NRB uncertainty statement should therefore avoid treating all terms as independent per-pixel random noise.

\subsubsection{Step 6: Product-level reporting}

For CANVAS, the uncertainty statement for NRB/ARD products should report at least:
\begin{itemize}
  \item inherited source-product radiometric and geometric uncertainty;
  \item uncertainty associated with DEM, local incidence angle, terrain correction, and terrain-normalisation convention;
  \item map-grid/geolocation uncertainty, including CRS, pixel coordinate convention, pixel spacing, and resampling method;
  \item mask-related uncertainty or quality flags for layover, shadow, no-data, and invalid pixels;
  \item correlation assumptions, including whether uncertainty terms are per-pixel, local-window, swath-scale, or product-wide;
  \item the traceability of ancillary data and processing versions used to generate the NRB product.
\end{itemize}

This subsection complements the SLC-focused uncertainty analysis above. It makes explicit that NRB/ARD procedures inherit source-acquisition uncertainty but also require a separate product-generation uncertainty budget for terrain correction, geocoding, masking, and metadata lineage.



